\section{Introduction}
\label{sec:introduction}

Classical simulation of quantum circuits is indispensable in the
current Noisy Intermediate-Scale Quantum (NISQ) era: it provides
the ground truth needed to validate quantum algorithms, characterise
noise models, and benchmark nascent hardware against classical
baselines. Yet full state-vector simulation is inherently expensive.
Representing the state of an $n$-qubit system requires storing
$2^n$ complex amplitudes, and at double precision this amounts to
over 16\,petabytes for 50~qubits~\cite{wang2021quantum}---far
exceeding the memory of any single machine and rivalling the
aggregate capacity of the largest supercomputers.

High-performance simulators have tackled this exponential barrier
through distributed memory, GPU acceleration, and careful gate
scheduling~\cite{wang2021quantum}, achieving impressive qubit counts
on dedicated HPC infrastructure. However, most of these systems
are designed for short-lived batch jobs: they provide limited or no
support for out-of-core execution on commodity hardware,
checkpointing during long-running simulations, or automatic
recovery after failure. When a 12-hour simulation crashes at
hour~11, the entire computation must typically be restarted from
scratch.

We argue that these gaps can be addressed by treating quantum
circuit simulation as a \emph{data management problem}. Database
systems have decades of experience managing datasets that exceed
main memory, ensuring transactional atomicity, and recovering
consistently after failures---techniques that map naturally onto the
lifecycle of a long-running state-vector simulation. Recent work by
Hai~et~al.~\cite{hai2025cqms} has begun to explore this connection,
demonstrating how relational modelling, structured logging, and
query-style interfaces can bring principled data management to
quantum computing workflows. Our work builds on this vision and
investigates the complementary question: can database-inspired
recovery and storage techniques enable fault-tolerant, large-scale
quantum state-vector simulation on commodity hardware?

We present an out-of-core state-vector simulator that separates
\emph{control} (scheduling, fault tolerance, metadata) from
\emph{data} (the numeric kernel operating on amplitudes). The
$2^n$-dimensional state is stored as dense binary chunks on disk and
streamed through CPU kernels, while Apache~Spark is retained purely
for task distribution across nodes---no amplitude data passes through
Spark's shuffle or aggregation machinery. Crash safety is provided
by double-buffered storage with a write-ahead log (WAL) that
guarantees at most one step of lost work per failure. Circuit
partitioning follows an Atlas-inspired staging
strategy~\cite{atlas} that minimises cross-chunk communication, and
empirical chunk-size tuning ensures the per-chunk working set remains
cache-friendly.

Our contributions are:
\begin{enumerate}
  \item An out-of-core streaming architecture for state-vector
    simulation that decouples the numeric kernel from the
    orchestration layer, achieving $300$--$10\,000\times$ speedup
    over a Spark~DataFrame baseline on identical circuits.
  \item A database-inspired recovery protocol---double-buffered
    storage with atomic writes and a step-level WAL---that enables
    crash-recoverable progress with bounded work loss.
  \item Integration of Atlas-style ILP circuit staging, single-qubit
    gate fusion, and level batching to reduce I/O passes and kernel
    invocations.
  \item An empirical evaluation demonstrating correctness across
    147~MQT~Bench circuit instances (up to 395\,059~gates),
    scaling to $n = 30$ ($>10^9$~amplitudes, 17\,GB on disk) on a
    single machine, and systematic chunk-size tuning that reveals
    a cache-driven performance cliff.
\end{enumerate}

\section{Related Work}
\label{sec:related}

This section reviews prior work relevant to scalable state-vector
simulation and positions our contribution at the intersection of
high-performance simulation and data management. We group related
efforts into five categories.

\paragraph{High-performance state-vector simulators.}
Full state-vector simulators such as those deployed on the Sunway
TaihuLight~\cite{wang2021quantum} and Summit supercomputers
distribute the $2^n$-element amplitude vector across thousands of
nodes, achieving 40--50~qubit simulations through careful MPI
communication patterns and GPU offloading. These systems maximise
throughput under the assumption of reliable infrastructure, but
typically offer no checkpointing or crash recovery: a single node
failure aborts the entire job.

\paragraph{Locality-aware scheduling and partitioning.}
The exponential state size makes communication the dominant cost in
distributed simulation. Atlas~\cite{atlas} formulates circuit
partitioning as an integer linear programme that assigns gates to
stages and selects local qubit sets to minimise inter-node shuffles.
Cache-blocking and gate fusion techniques further reduce memory
traffic within each node. Our implementation adopts Atlas's ILP
staging and extends it with insular-qubit detection for diagonal
gates.

\paragraph{Big-data frameworks for numerical workloads.}
Apache Spark and similar dataflow engines provide built-in fault
tolerance (lineage-based recomputation) and distributed scheduling.
Prior work, including earlier versions of this system (v1--v3,
Section~\ref{sec:prior-versions}), has modelled amplitudes as
DataFrame rows and applied gates via join-aggregate queries.
While functionally correct, the per-operator overhead of Spark's
execution engine (job scheduling, serialisation, shuffle I/O)
dominates the cost of the small $2\!\times\!2$ matrix updates,
making this approach orders of magnitude slower than direct
computation.

\paragraph{Out-of-core numerical methods.}
External-memory algorithms for linear algebra---blocked matrix
multiply, out-of-core LU, and streaming SVD---have a long history
in scientific computing. Our chunk-based streaming approach applies
the same principle to state-vector simulation: the amplitude vector
is partitioned into cache-sized blocks, streamed sequentially
through the kernel, and written back to disk, bounding memory
consumption independently of~$n$.

\paragraph{Database-inspired fault tolerance.}
Write-ahead logging, atomic commit protocols, and double-buffering
are foundational techniques in database recovery~\cite{mohan1992aries}.
Hai~et~al.~\cite{hai2025cqms} recently advocated applying such
data management principles to quantum computing workflows, including
simulation metadata, provenance tracking, and structured logging.
Our system operationalises this vision for the simulation kernel
itself: every step of state evolution is protected by a WAL and
atomic buffer swap, ensuring consistent recovery after fail-stop
failures with at most one step of repeated work.


%% ----------------------------------------------------------------
%% Implementation section follows
%% ----------------------------------------------------------------

\section{Implementation}
\label{sec:implementation}

We developed four successive prototype versions of a quantum state-vector simulator.
Versions v1--v3 explored relational formulations where amplitudes are stored as
database rows and gate application is expressed as join-aggregate queries. Version~v4
(the current system) abandons the relational data plane entirely in favour of an
out-of-core streaming kernel, while retaining Spark for orchestration and fault
tolerance. We briefly summarise v1--v3 before describing v4 in detail.

\subsection{Prior Versions (v1--v3)}
\label{sec:prior-versions}

\paragraph{v1: SQLite baseline.}
The initial prototype stored amplitudes in a SQLite table
\texttt{(version, idx, real, imag)} and implemented gate application as SQL
join-and-aggregate queries over the unitary matrix entries.
Each gate was executed inside an explicit transaction with a WAL entry
(\texttt{PENDING} $\to$ \texttt{COMMITTED}) to guarantee atomic progress.
This version served as a correctness reference and demonstrated that the
relational formulation is functionally correct, but performance was limited
by per-gate transaction overhead.

\paragraph{v2: Spark DataFrames.}
v2 ported the relational model to Apache Spark. Amplitudes were represented as
a distributed DataFrame \texttt{(idx, real, imag)}, persisted to Parquet between
batches. Gate application used the same join-aggregate semantics implemented with
Spark primitives (broadcast join, \texttt{groupBy}, complex multiply-add). A
DuckDB-backed metadata store tracked WAL entries and checkpoints. Recovery on
restart resumed from the last committed Parquet state version.

\paragraph{v3: topological levels and tensor-product fusion.}
v3 introduced a dependency DAG over gates, grouping qubit-disjoint gates into
topological levels. Within each level, independent single-qubit gates were fused
via their tensor product into a single broadcast-join-aggregate pass, reducing the
number of full-state transformations. Two-qubit gates were applied sequentially.
While this improved throughput for circuits with large layers of independent
single-qubit gates, every gate application still paid the cost of Spark job
scheduling ($\sim$200--500\,ms), serialisation, and shuffle I/O.
Benchmarking showed that v3's overhead made it 300--10\,000$\times$ slower than
direct NumPy computation for circuits up to $n=12$ (Table~\ref{tab:v3-comparison}).

\subsection{v4: Out-of-Core Streaming Kernel}
\label{sec:v4}

Guided by the observation that the core operation is a small
$2\!\times\!2$ or $4\!\times\!4$ unitary update applied repeatedly
across a long amplitude vector, v4 restructures the simulator to
separate the \emph{control plane} (scheduling, fault tolerance,
metadata) from the \emph{data plane} (numeric kernel).
The state vector $\psi \in \mathbb{C}^{2^n}$ is stored on disk as
dense chunk files and streamed through CPU kernels one chunk at a time.
Spark is retained purely for task distribution across nodes; no
DataFrame joins, shuffles, or aggregations touch the amplitude data.

\subsubsection{State Storage}
\label{sec:storage}

The $2^n$ amplitudes are stored as \texttt{complex64} (single-precision
complex, 8 bytes per amplitude) in contiguous binary chunk files.
Each chunk holds $C = 2^k$ amplitudes for a configurable $k$; chunk~$i$
covers the index range $[i \cdot C,\;(i{+}1) \cdot C)$. A JSON manifest
in each buffer directory records the number of qubits, chunk size, file
list, and creation timestamp. Every file write follows an atomic
protocol: write to a temporary file, call \texttt{fsync}, then
\texttt{os.replace} to the final path---ensuring that a crash at any
point leaves either the old or the new file intact.

Two buffer directories (\texttt{state\_a}, \texttt{state\_b}) alternate
as source and destination (double-buffering). The source buffer is
\emph{never modified} during a step; only the destination is written.
This invariant guarantees that the source is always a valid, consistent
snapshot of the state, enabling safe crash recovery without
application-level undo.

\subsubsection{Local and Non-Local Gate Classification}
\label{sec:local-nonlocal}

Given chunk size $C = 2^k$, the lowest $k$ bits of the amplitude index
identify the position \emph{within} a chunk, while the upper $n{-}k$
bits identify \emph{which chunk}. A gate on qubit~$q$ pairs amplitudes
differing only in bit~$q$:
\begin{itemize}
  \item \textbf{Local} ($q < k$): both amplitudes reside in the same
    chunk. The gate is applied by reading a single chunk, performing the
    update in-place, and writing the result.
  \item \textbf{Non-local} ($q \geq k$): the paired amplitudes lie in
    different chunks. We load a \emph{partner group} of chunks (those
    differing only in bit~$q$) and apply the gate across the group using
    a butterfly exchange pattern~\cite{fang2022hisvsim}.
\end{itemize}
For two-qubit gates, four cases arise depending on whether each target
qubit is local or non-local, requiring pairs or quads of chunks.

\subsubsection{Kernels}
\label{sec:kernels}

All kernels operate on in-memory NumPy arrays and perform no I/O.
Two CPU kernel implementations are provided:

\paragraph{Scalar kernel.}
For each pair (1Q) or quad (2Q) of index positions within a chunk,
the kernel applies the $2\!\times\!2$ or $4\!\times\!4$ matrix update
via explicit element-wise arithmetic. This kernel is simple and easy to
verify.

\paragraph{Batched GEMM kernel.}
Index pairs are gathered into a $(2 \times M)$ or $(4 \times M)$ matrix
(where $M = C/2$ or $C/4$), multiplied by the gate matrix via NumPy's
BLAS-backed \texttt{matmul}, and scattered back. This leverages
optimised BLAS routines and is faster when the working set fits in
CPU cache.

\paragraph{Non-local kernel.}
For non-local gates, pure array operations combine partner chunks:
element-wise $2\!\times\!2$ for single-qubit non-local gates, and
element-wise $4\!\times\!4$ for fully non-local two-qubit gates.
Mixed cases (one qubit local, one non-local) combine local pair
extraction with cross-chunk element-wise operations.

\subsubsection{Circuit Optimisations}
\label{sec:circuit-opt}

\paragraph{Levelisation and level batching.}
Gates are grouped into dependency-free \emph{levels} via a topological
sort on the per-qubit ordering DAG. Consecutive levels containing only
local gates are merged into a single I/O step (one read--compute--write
pass over all chunks instead of one per level).

\paragraph{Single-qubit gate fusion.}
Consecutive single-qubit gates on the same qubit (with no intervening
two-qubit gate on that qubit) are pre-multiplied into a single
$2\!\times\!2$ matrix before execution, reducing the number of kernel
invocations.

\paragraph{Atlas-style circuit staging.}
For circuits where many gates target qubits in the non-local range
($q \geq k$), we implement circuit staging inspired by Atlas~\cite{atlas}.
The circuit is partitioned into \emph{stages}, each associated with a
set of $k$ \emph{local qubits}. Between stages, SWAP operations
physically rearrange the state vector so that the next stage's target
qubits occupy positions $0\ldots k{-}1$.

Three partitioning methods are provided:
\begin{enumerate}
  \item \textbf{ILP} (requires PuLP): an integer linear programme that
    assigns gates to stages and selects local qubit sets to minimise
    the total number of qubit transitions, faithfully reimplemented from
    Atlas~\cite{atlas}.
  \item \textbf{Heuristic}: a dependency-aware greedy algorithm ported
    from Atlas's reference implementation, iteratively executing gates
    whose non-insular qubits are in the current local set and selecting
    the next local set by priority (qubits in the first blocked gate
    $>$ global gate count $>$ local gate count).
  \item \textbf{Greedy}: a simpler frequency-based lookahead from our
    earlier prototype.
\end{enumerate}
Following Atlas, sparse (diagonal) gates---Z, S, T, CZ, CR---are
treated as \emph{insular}: their qubits need not be local, as the
conditional-phase update can be applied without a full shuffle. This
reduces the number of required stages.

\subsubsection{Runners}
\label{sec:runners}

\paragraph{Sequential runner.}
Processes work items (chunk groups) one at a time. Memory footprint is
minimal: at most one chunk group resides in RAM at a time.

\paragraph{Pipeline runner.}
Three threads---reader, compute worker, writer---connected by bounded
queues of configurable depth. This overlaps disk I/O with computation.
Non-local gate groups are processed outside the pipeline (they require
multiple chunks loaded simultaneously).

\paragraph{Spark runner.}
Spark parallelises local-gate chunk processing across executors. Each
executor reads its assigned chunk from disk, applies the kernel, and
writes the result. Non-local groups are processed on the driver.
Spark is used \emph{only for task distribution}---no amplitude data
passes through Spark's shuffle or aggregation machinery.

\subsubsection{Write-Ahead Log and Crash Recovery}
\label{sec:wal}

A single atomic JSON file (\texttt{wal.json}) records the circuit hash,
which buffer (\texttt{a} or \texttt{b}) holds the latest committed
state, and how many steps have been completed. Each step follows the
protocol:
\begin{enumerate}
  \item Read chunks from the committed source buffer.
  \item Apply the kernel; write results to the destination buffer
    (atomic per-chunk writes).
  \item Write the destination manifest atomically.
  \item Update the WAL to flip the committed buffer and increment the
    step counter (atomic write).
\end{enumerate}
On crash, recovery reads the WAL, identifies the committed buffer as
the valid source, wipes the destination buffer (which may contain
partial data), and resumes from the recorded step. Because the source
buffer is never modified, the input to any interrupted step is always
intact, guaranteeing exactly-once step semantics under fail-stop
failures. A fencing lock file prevents concurrent processes from
operating on the same working directory.

\section{Experimental Evaluation}
\label{sec:results}

\subsection{Experimental Setup}
All experiments were conducted on a MacBook Pro with Apple M3 Pro
(14~cores), 48\,GB unified memory, and 460\,GB NVMe SSD
(${\sim}3.5$\,GB/s sequential read). The simulator uses Python~3.9
with NumPy (BLAS-backed) for kernels and PySpark in
\texttt{local[*]} mode for the Spark runner.

\subsection{Correctness Validation}
\label{sec:correctness}

We validate correctness by running each circuit through our simulator
and comparing the resulting state vector against two independent
oracles:
\begin{enumerate}
  \item \textbf{ref\_dense}: an in-memory reference simulator that
    applies gates one-by-one to a NumPy state vector (complex128).
    Used as ground truth for kernel-level and end-to-end tests.
  \item \textbf{Qiskit Statevector}: Qiskit's built-in statevector
    simulator, a mature and widely trusted implementation.
\end{enumerate}
To ensure coverage beyond hand-written test circuits, we draw
147~circuit instances from 30~families in the MQT~Bench
suite~\cite{mqtbench}---a standardised collection of quantum
algorithm circuits including GHZ, QFT, Grover, Shor, random circuits,
variational ans\"{a}tze, and arithmetic circuits. Each MQT~Bench
circuit is transpiled to our basis gate set and then executed by
(i)~our out-of-core runner, (ii)~ref\_dense, and (iii)~Qiskit
Statevector. We compare the three output state vectors; all
147~instances for $n = 3$ to $20$ qubits (with gate counts ranging
from 3 for GHZ-3 to 395\,059 for Shor-18) achieve overlap
$|\langle\psi_{\text{oracle}}|\psi_{\text{ours}}\rangle| > 1 -
10^{-6}$ against both oracles. Little-endian qubit ordering is
enforced by a dedicated test: applying $X$ to qubit~0 on
$|0\rangle^n$ must produce amplitude $1.0$ at index~1.

The full test suite comprises 137 unit tests covering circuit
validation, endianness, kernel correctness (scalar, batched, non-local),
out-of-core end-to-end vs.\ oracle, crash recovery (12 tests with
simulated crashes via subprocesses), WAL integrity, double-buffer
alternation, fencing lock semantics, Spark runner, gate fusion, and
Atlas-style staging with qubit permutation.

\subsection{Comparison with v3}
\label{sec:v3-comparison}

Table~\ref{tab:v3-comparison} compares v3 (Spark DataFrame gate
application) against v4 (out-of-core kernel) on identical circuits.
Correctness is verified: all state vectors match (overlap $= 1.0$).

\begin{table}[h]
\centering
\caption{Runtime comparison between v3 (Spark) and v4 (kernel).}
\label{tab:v3-comparison}
\begin{tabular}{lrrrrr}
\toprule
Circuit & $n$ & Gates & v3 (s) & v4 (s) & Speedup \\
\midrule
GHZ-3            &  3 &     3 &   3.64 & 0.003 & 1\,097$\times$ \\
GHZ-10           & 10 &    10 &   2.65 & 0.009 &    310$\times$ \\
GHZ-12           & 12 &    12 &   2.98 & 0.011 &    278$\times$ \\
QFT-6            &  6 &    21 &   3.51 & 0.008 &    417$\times$ \\
QFT-8            &  8 &    36 &   5.52 & 0.015 &    362$\times$ \\
MQT-dj-10        & 10 &    62 &  14.82 & 0.025 &    585$\times$ \\
MQT-graphstate-10& 10 &    20 &   5.80 & 0.004 & 1\,318$\times$ \\
H-wall-4         &  4 &     4 &   2.80 & 0.002 & 1\,356$\times$ \\
H-wall-8         &  8 &     8 &   0.86 & 0.002 &    418$\times$ \\
H-wall-12        & 12 &    12 &  94.87 & 0.009 &10\,087$\times$ \\
\bottomrule
\end{tabular}
\end{table}

The median speedup is ${\sim}320\times$, with extremes exceeding
$10\,000\times$. The speedup is attributable to eliminating Spark job
scheduling, serialisation, and shuffle overhead from the critical path.

\subsection{Scaling: Non-Stabilizer Circuits}
\label{sec:scaling}

To characterise scaling limits, we use non-stabilizer circuits
(H~+~T~+~CNOT layers) that produce fully dense state vectors, ensuring
the exponential state size is the binding constraint. All results
include norm verification ($\|\psi\|_2 = 1 \pm 10^{-4}$).

\paragraph{In-memory (ref\_dense, complex128).}
Table~\ref{tab:scaling-inmem} shows that the in-memory simulator
handles up to $n = 28$ (4.3\,GB state, 23\,GB peak RSS) in 6.9
minutes on this machine. At $n = 29$, estimated peak memory
(${\sim}46$\,GB) exceeds physical RAM.

\begin{table}[h]
\centering
\caption{In-memory scaling (ref\_dense, complex128, non-stabilizer circuits).}
\label{tab:scaling-inmem}
\begin{tabular}{rrrrl}
\toprule
$n$ & Amplitudes & State & Time & Peak RAM \\
\midrule
20 &        1\,M &   17\,MB &    1.1\,s & 141\,MB \\
24 &       16\,M &  268\,MB &   15.5\,s & 1.8\,GB \\
26 &       67\,M & 1.07\,GB &   68.6\,s & 6.9\,GB \\
28 &      268\,M & 4.30\,GB &  412\,s   & 23\,GB \\
\bottomrule
\end{tabular}
\end{table}

\paragraph{Out-of-core (on-disk, complex64).}
Table~\ref{tab:scaling-ooc} demonstrates that the out-of-core runner
extends the frontier to $n = 30$ (over 1 billion amplitudes, 17\,GB
on disk, 1\,024 chunks) in 10.5 minutes. The runner's memory footprint
is bounded by a constant number of chunks in RAM, independent of~$n$.
Correctness is cross-validated against ref\_dense at $n = 20$.

\begin{table}[h]
\centering
\caption{Out-of-core scaling (on-disk, complex64, non-stabilizer circuits).}
\label{tab:scaling-ooc}
\begin{tabular}{rrrrl}
\toprule
$n$ & Amplitudes & Disk & Chunks & Time \\
\midrule
20 &        1\,M &  0.02\,GB &      1 &    1.2\,s \\
24 &       16\,M &  0.27\,GB &     16 &   15.4\,s \\
26 &       67\,M &  1.07\,GB &     64 &   61\,s \\
28 &      268\,M &  4.29\,GB &    256 &  139\,s \\
29 &      537\,M &  8.59\,GB &    512 &  285\,s \\
30 &    1\,074\,M & 17.2\,GB & 1\,024 &  631\,s \\
\bottomrule
\end{tabular}
\end{table}

\subsection{Chunk Size Tuning}
\label{sec:chunk-tuning}

The chunk size $C = 2^k$ is the primary hyperparameter: it determines
how many qubits are local ($k$), the number of I/O passes, and whether
the per-chunk working set fits in CPU cache. We swept chunk sizes
$2^{14}$ through $2^{22}$ on non-stabilizer circuits at $n = 24$ and
$n = 26$, testing both the sequential and pipeline runners with and
without gate fusion.

\begin{table}[h]
\centering
\caption{Chunk size sweep at $n = 26$ (537\,MB state, 154 gates,
pipeline runner, fusion enabled, buffer depth 4).}
\label{tab:chunk-sweep}
\begin{tabular}{rrrrr}
\toprule
$k$ & Chunk size & Chunks & Steps & Time (s) \\
\midrule
14 &   0.1\,MB & 4\,096 & 19 & 63.5 \\
16 &   0.5\,MB & 1\,024 & 17 & \textbf{40.2} \\
18 &   2.1\,MB &    256 & 15 & \textbf{40.0} \\
20 &   8.4\,MB &     64 & 13 & 118.9 \\
22 &  33.6\,MB &     16 & 11 & 116.9 \\
\bottomrule
\end{tabular}
\end{table}

The results reveal a non-monotonic relationship between chunk size and
wall time (Table~\ref{tab:chunk-sweep}). Although larger chunks reduce
the number of I/O steps (from 19 at $k{=}14$ to 11 at $k{=}22$),
the optimal wall time is achieved at $k = 16$--$18$ (0.5--2\,MB
chunks). Beyond $k = 18$, the gather/scatter kernel working set
exceeds the CPU's L2 cache (${\sim}4$\,MB per core on Apple Silicon),
causing a ${\sim}3\times$ slowdown that outweighs the reduction in
I/O passes. Gate fusion consistently provides a 10--40\% speedup
by merging consecutive local levels and pre-multiplying single-qubit
gate matrices. The pipeline runner offers a 10--15\% improvement over
the sequential runner at the optimal chunk size.

\subsection{Crash Recovery}
\label{sec:recovery-eval}

We evaluate recovery correctness under simulated fail-stop failures.
Crashes are injected by terminating the runner subprocess
(via \texttt{os.\_exit}) at controlled points: after initialisation,
after writing a subset of chunks, and after completing a full step
but before WAL commit. In each case, the recovery procedure reads
the WAL, identifies the committed buffer, wipes the (potentially
corrupted) destination buffer, and resumes from the last committed
step. Across 12 dedicated recovery tests, the final state after
recovery matches the fault-free reference in all cases. At most
one step of work is lost per crash, bounded by the step-level WAL
granularity.

\subsection{Summary}
\label{sec:results-summary}

The current implementation (v4) demonstrates:
(i)~correctness against two independent oracles across 147 circuit
instances spanning 3--20 qubits and up to 395\,059 gates;
(ii)~a $300$--$10\,000\times$ speedup over the Spark DataFrame
approach (v3) by moving the numeric kernel out of the Spark execution
path;
(iii)~out-of-core scaling to $n = 30$ (over $10^9$ amplitudes, 17\,GB
on disk) on a single machine with bounded memory footprint;
(iv)~empirically tuned chunk size ($2^{16}$--$2^{18}$) that balances
cache efficiency against I/O pass count; and
(v)~atomic, crash-recoverable progress via double-buffered storage
with a step-level write-ahead log.

\bibliographystyle{plain}
\begin{thebibliography}{9}

\bibitem{hai2025cqms}
Rihan Hai, Shih-Han Hung, Tim Coopmans, Tim Littau, and Floris Geerts.\\
\newblock Quantum Data Management in the NISQ Era.\\
\newblock \emph{Proceedings of the VLDB Endowment}, 18(6):1720--1729, February 2025.\\
\newblock \href{https://doi.org/10.14778/3725688.3725701}{doi:10.14778/3725688.3725701}.

\bibitem{wang2021quantum}
Zhimin Wang, Zhaoyun Chen, Shengbin Wang, Wendong Li, Yongjian Gu, Guoping Guo, and Zhiqiang Wei.\\
\newblock A quantum circuit simulator and its applications on Sunway TaihuLight supercomputer.\\
\newblock \emph{Scientific Reports}, 11(1):355, 2021.\\
\newblock \href{https://doi.org/10.1038/s41598-020-79777-y}{doi:10.1038/s41598-020-79777-y}.

\bibitem{jones2019quest}
Tyson Jones, Anna Brown, Ian Bush, and Simon C. Benjamin.
\newblock QuEST and high performance simulation of quantum computers.
\newblock \emph{Scientific Reports}, 9(1):10736, 2019.
\newblock \href{https://doi.org/10.1038/s41598-019-47174-9}{doi:10.1038/s41598-019-47174-9}.

\bibitem{guerreschi2020intel}
Gian Giacomo Guerreschi, Justin Hogaboam, Fabio Baruffa, and Nicolas P. D. Sawaya.
\newblock Intel Quantum Simulator: A cloud-ready high-performance simulator of quantum circuits.
\newblock \emph{Quantum Science and Technology}, 5(3):034007, 2020.
\newblock \href{https://doi.org/10.1088/2058-9565/ab8505}{doi:10.1088/2058-9565/ab8505}.

\bibitem{fang2022hisvsim}
Bo Fang, M. Yusuf \"{O}zkaya, Ang Li, \"{U}mit V. \c{C}ataly\"{u}rek, and Sriram Krishnamoorthy.
\newblock Efficient Hierarchical State Vector Simulation of Quantum Circuits via Acyclic Graph Partitioning.
\newblock In \emph{2022 IEEE International Conference on Cluster Computing (CLUSTER)}, pages 289--300. IEEE, 2022.
\newblock \href{https://doi.org/10.1109/CLUSTER51413.2022.00041}{doi:10.1109/CLUSTER51413.2022.00041}.

\bibitem{cambiucci2024hypergraphic}
Waldemir Cambiucci, Regina Melo Silveira, and Wilson Vicente Ruggiero.
\newblock Hypergraphic partitioning for spatial and temporal quantum circuit cutting.
\newblock In \emph{2024 IEEE International Conference on Quantum Computing and Engineering (QCE)}, pages 486--487. IEEE, 2024.
\newblock \href{https://doi.org/10.1109/QCE60285.2024.10368}{doi:10.1109/QCE60285.2024.10368}.

\bibitem{gu2015marlin}
Rong Gu, Yin Hu, and Yihua Huang.
\newblock Marlin: A fast and scalable distributed matrix multiplication library on Spark.
\newblock In \emph{2015 IEEE International Conference on Big Data (Big Data)}, pages 2539--2552. IEEE, 2015.
\newblock \href{https://doi.org/10.1109/BigData.2015.7364023}{doi:10.1109/BigData.2015.7364023}.

\bibitem{gittens2016matrix}
Alex Gittens et al.
\newblock Matrix factorizations at scale: a comparison of scientific data analytics in Spark and C+MPI using three case studies.
\newblock In \emph{2016 IEEE International Conference on Big Data (Big Data)}, pages 204--213. IEEE, 2016.
\newblock \href{https://doi.org/10.1109/BigData.2016.7840606}{doi:10.1109/BigData.2016.7840606}.

\bibitem{sun2025accelerating}
Wenbo Sun, Asterios Katsifodimos, and Rihan Hai.
\newblock Accelerating machine learning queries with linear algebra query processing.
\newblock \emph{Distributed and Parallel Databases}, 43:8, 2025.
\newblock \href{https://doi.org/10.1007/s10619-024-07451-7}{doi:10.1007/s10619-024-07451-7}.

\bibitem{gonthier2025scheduler}
Maxime Gonthier, Loris Marchal, and Samuel Thibault.
\newblock A scheduler to foster data locality for GPU and out-of-core task-based linear algebra applications.
\newblock \emph{Journal of Parallel and Distributed Computing}, 206:105170, August 2025.
\newblock \href{https://doi.org/10.1016/j.jpdc.2025.105170}{doi:10.1016/j.jpdc.2025.105170}.

\bibitem{toledo1999survey}
Sivan Toledo.
\newblock A survey of out-of-core algorithms in numerical linear algebra.
\newblock In \emph{External Memory Algorithms}, volume 50 of \emph{DIMACS Series in Discrete Mathematics and Theoretical Computer Science}, pages 161--179. American Mathematical Society, 1999.
\newblock \href{https://doi.org/10.1090/dimacs/050}{doi:10.1090/dimacs/050}.

\bibitem{mohan1992aries}
C. Mohan, Don Haderle, Bruce Lindsay, Hamid Pirahesh, and Peter Schwarz.
\newblock ARIES: a transaction recovery method supporting fine-granularity locking and partial rollbacks using write-ahead logging.
\newblock \emph{ACM Transactions on Database Systems (TODS)}, 17(1):94--162, 1992.
\newblock \href{https://doi.org/10.1145/128765.128770}{doi:10.1145/128765.128770}.

\bibitem{atlas}
Mingkuan Xu, Shiyi Cao, Xupeng Miao, Umut A. Acar, and Zhihao Jia.\\
\newblock Atlas: Hierarchical Partitioning for Quantum Circuit Simulation on GPUs (Extended Version).\\
\newblock In \emph{Proceedings of the International Conference for High Performance Computing, Networking, Storage and Analysis (SC '24)}, 2024.\\
\newblock \href{https://doi.org/10.48550/arXiv.2408.09055}{doi:10.48550/arXiv.2408.09055}.

\bibitem{mqtbench}
Nils Quetschlich, Lukas Burgholzer, and Robert Wille.
\newblock MQT Bench: Benchmarking Software and Design Automation Tools for Quantum Computing.
\newblock \emph{Quantum}, 7:1062, 2023.
\newblock \href{https://doi.org/10.22331/q-2023-07-20-1062}{doi:10.22331/q-2023-07-20-1062}.

\end{thebibliography}
